\chapter{Lab 2}


\section{评分标准}
\textbf{60\%验收+40\%报告}

验收部分十个题,每个同学问两个总共占验收分数的30\%，70\% 实验部分由 50\% 仿真和 20\% 上板组成。


\section{预设问题}

\begin{enumerate}
    \item \textbf{Q：}实验中 \texttt{if\_stall} 和 \texttt{mem\_stall} 这两个信号的作用有什么不同?

    \textbf{A：}
    \begin{itemize}
        \item \texttt{if\_stall}: 主要用于取指阶段。当正在等待指令从内存返回时，它会暂停程序计数器（PC）的更新，确保在当前指令成功取回前，不会去取下一条指令。
        \item \texttt{mem\_stall}: 主要用于处理访存冲突。当一条访存指令（如 \texttt{ld/sw}）在 MEM 阶段需要访问总线，而总线正被取指占用时，它会暂停流水线的前几个阶段（PC, IF/ID, ID/EXE, EXE/MEM），让访存请求可以等待并最终完成，同时清空（flush）后续无效的流水段。
    \end{itemize}

    \item \textbf{Q：}为什么在本实验的架构中，CPU 无法同时进行取指令和数据访存？这属于哪一种类型的竞争？

    \textbf{A：}因为指令内存和数据内存在物理上是同一个 AXI BRAM 模块，它们共享同一个总线接口。当流水线中的一条指令在取指阶段需要读内存，而另一条指令在访存阶段也需要读/写内存时，它们会争用同一个物理资源。这属于结构竞争。

    \item \textbf{Q：}为什么要设计一个取指/访存状态机来与 AXI 总线交互，而不是像 Lab1 一样直接读写内存？

    \textbf{A：}因为真实的内存访问不是单周期完成的，它需要经过请求、等待、响应等多个步骤，耗时是不确定的。设计状态机是为了管理这个复杂的、多周期的握手过程，并处理取指和访存同时发生时的仲裁问题，确保每次只有一个请求在总线上被正确处理。

    \item \textbf{Q：}当 ID 阶段的指令需要暂停（Stall）时，流水线的前级（IF/ID）和后级（EX/MEM/WB）分别应该做什么操作？

    \textbf{A：}
    \begin{itemize}
        \item 前级 (IF, ID): 应当被“锁住”或“冻结”，即保持当前状态，不接收新的指令，PC 也不更新。
        \item 后级 (EX, MEM, WB): 应当继续正常执行，或者被插入无效指令，让已经在流水线后段的指令能够继续流动并最终完成，从而为被暂停的指令“腾出空间”。
    \end{itemize}

    \item \textbf{Q：}在流水线中，一个操作数可能同时被多个后续阶段的指令写入。如果 ID 阶段的一条指令同时检测到来自 EXE 阶段和 MEM 阶段的 forwarding 路径都有效，它应该优先选择哪一个？为什么？

    \textbf{A：}它应该优先选择来自 EXE 阶段的转发数据。原因是：流水线遵循程序的顺序执行语义。来自 EXE 阶段的数据是由更“新”（即更靠后）的指令产生的。选择来自最靠近 ID 阶段（即 EXE 阶段）的数据，可以确保指令获取到的是最新版本的值，这与指令在程序中的顺序逻辑一致。如果错误地选择了来自 MEM 阶段的数据，将会使用一个已经被后续指令覆盖的“陈旧”值。

    \item \textbf{Q：}为什么当 \texttt{mem\_stall} 信号有效时，需要暂停 PC、IF/ID、ID/EXE、EXE/MEM 寄存器，同时却要 flush MEM/WB 寄存器？

    \textbf{A：}暂停前级流水段是为了“冻结”流水线，让访存指令停留在 MEM 阶段，等待总线资源，并确保后续指令不会错误地进入执行阶段。而清空 MEM/WB 寄存器是为了防止访存指令在 stall 的第一个周期就错误地将无效数据写入写回段，因为访存操作尚未完成，还没有有效的数据可以传递给 WB 阶段。这确保了流水线后级可以排空，为暂停的指令腾出空间。

    \item \textbf{Q：}\texttt{if\_stall} 信号在状态机中的主要作用是什么？在哪些状态下它会被激活，为什么？

    \textbf{A：}\texttt{if\_stall} 的主要作用是暂停 PC 的更新，确保在当前指令成功从内存取回之前，CPU 不会去取下一条指令。它在所有涉及等待内存响应的状态下都会被激活，包括 IF1, IF2, WAITFOR1, WAITFOR2, MEM1, MEM2。

    \item \textbf{Q：}如果将 AxiBram 替换为一个响应速度的内存模型（例如，需要几十个周期才能返回数据），当前的状态机设计是否还能正常工作？

    \textbf{A：}状态机设计本身仍然可以正常工作。因为它是基于 \texttt{valid/ready} 来驱动状态转换的，无论设备响应多慢，状态机都会停留在等待状态直到收到 \texttt{reply\_valid} 信号。

    \item \textbf{Q：}在实验的取指/访存状态机中，为什么要引入 \texttt{WAITFOR1} 和 \texttt{WAITFOR2} 这两个状态？它们解决了什么问题？

    \textbf{A：}这两个状态是为了解决取指和访存同时请求总线时的结构竞争问题。当 MEM 阶段的访存指令需要访问总线，而 IF 阶段的取指请求已经发出时，总线必须先完成当前的取指操作。\texttt{WAITFOR} 状态会暂停流水线，等待取指操作完成，然后再将总线控制权交给访存操作，进入 MEM1 状态。

    \item \textbf{Q：}当流水线因为 load-use 数据冒险（例如 \texttt{ld x1, 0(sp)} 后紧跟 \texttt{add x2, x1, x3}）而暂停时，最少需要暂停几个周期？

    \textbf{A：}1个。Forwarding 可以解决剩余延迟，具体流程如下：
    \begin{enumerate}
        \item \texttt{ld} 指令在 MEM 阶段从内存中取回数据。
        \item 当 \texttt{add} 指令位于 ID 阶段时，它需要的 x1 值还不可用，因此流水线必须暂停一个周期。
        \item 在这个暂停周期中，\texttt{ld} 指令从 MEM 阶段前进到 WB 阶段，而 \texttt{add} 指令则停留在 ID 阶段。
        \item 暂停结束后，\texttt{add} 指令进入 EXE 阶段。此时，\texttt{ld} 指令的结果已经存在于 MEM/WB 流水线寄存器中。这个值可以通过 MEM -> EXE forwarding 直接传送给 \texttt{add} 指令的 ALU 输入端，使其能够正确执行。因此，一个周期的暂停就足以弥补访存操作带来的一个周期延迟。
    \end{enumerate}
\end{enumerate}